\documentclass[12pt]{article}
\usepackage[utf8]{inputenc}
\usepackage[T1]{fontenc}
\usepackage{hyperref}
\title{Pawaia Language Resources for K-12 Curriculum Development}
\author{Compiled by Claude Sonnet 5.5 (claude-sonnet-5-5) for Andr\'as Kornai}
\date{}
\begin{document}
\maketitle

\section{Introduction}
Pawaia (ISO 639-3 \texttt{pwa}; also Aurama, Sira, Tudahwe, Yasa) is spoken along the Pio and Purari rivers in Gulf and Simbu provinces, Papua New Guinea. Speaker estimates in search results range from 2,000--2,600 (1972) to 6,000 (2000). Its classification is unresolved: Ross treats it as a tentative Trans-New Guinea branch, while others leave it unclassified. No dedicated dictionary, primer or school programme was found.

\section{Primary Educational Resources}
No primers, textbooks or school materials were found in search results.

\section{Linguistic and Reference Materials}
\begin{itemize}
\item Wikipedia, Pawaia language. \url{https://en.wikipedia.org/wiki/Pawaia_language}
\item New Guinea World, Pawaia (Tua River family): overview and references. It cites D.\ Trefry, \emph{A Comparative Study of Kuman and Pawaian} (1969), which includes grammar notes on North Pawaia. \url{https://newguineaworld.linguistik.uzh.ch/families/papuan-gulf/tua-river/pawaia}
\item Trans New Guinea database, Pawaia. \url{https://transnewguinea.org/language/pawaia?page=2}
\item WALS entry for Pawaia. \url{https://wals.info/languoid/lect/wals_code_paw}
\item Endangered Languages Project, Pawaia bibliography. \url{https://www.endangeredlanguages.com/lang/10740/bibliography}
\end{itemize}

\section{Religious and Community Resources}
None were found in search results.

\section{Digital and Audio Resources}
\begin{itemize}
\item Wikitongues page for Pawaia (language profile; I did not verify whether it holds any recordings). \url{https://wikitongues.org/languages/pwa/}
\end{itemize}

\section{Curriculum Development Context}
Published description is limited to 1960s--70s work (Trefry, MacDonald), and no teaching materials were found. A curriculum would need fresh documentation with the community, and could build on PNG vernacular-education practice.

\end{document}